\section{Birational preparation}
\label{sec:preparation}

The purpose of this section is to replace the hypothetical case of maximal
nef dimension by a smooth threefold carrying a single connected nef cycle.
The torsion assertion below concerns the entire Cartier cycle, with its
nilpotents.

\begin{proposition}[The prepared cycle]
\label{prop:prepared-cycle}
Suppose that \(k\) is uncountable and algebraically closed of characteristic
\(p>3\), and that \(X\) is a projective terminal threefold with \(K_X\) nef,
\(\nu(K_X)=1\), and \(n(X,K_X)=3\).
There exist a smooth projective threefold \(V\), a generically finite
separable morphism \(\rho:V\to X\), a very ample divisor \(H\) on \(V\), and
a nonzero effective Cartier divisor
\[
 D=\sum_{\alpha}m_\alpha D_\alpha,\qquad
 m_\alpha\in\Z_{>0},\qquad \gcd_\alpha m_\alpha=1,
\]
with the following properties.
\begin{enumerate}[label=\textup{(\roman*)}]
\item The support of \(D\) is connected and has simple normal crossings.
The divisor \(D\) is nef, and \(D|_{D_\alpha}\) is numerically trivial
for every \(\alpha\).
\item There is no \(D\)-trivial curve through a very general point of \(V\).
Moreover \(D\cdot H^2>0\) and \(D^2\cdot H=0\).
\item The invertible sheaf \(\cO_D(D)\) is torsion.  Its order will be denoted
by \(r\).
\item For an integer \(b\geq0\), an open neighborhood \(U\) of \(|D|\), and
integers \(b_\alpha\), there is an isomorphism
\begin{equation}
\label{eq:prepared-canonical}
 \omega_V^{\otimes p^b}|_U
 \simeq \cO_U(B_*),\qquad B_*=\sum_\alpha b_\alpha D_\alpha .
\end{equation}
Consequently \(K_V\cdot D\cdot A=0\) for every Cartier divisor \(A\) on \(V\).
\end{enumerate}
\end{proposition}

The case $n(X,K_X)\leq2$ is already covered by
\cite[Theorem~1.8]{xu2024}, so this is the remaining nef-dimension case.
The rest of this section proves Proposition~\ref{prop:prepared-cycle}.
Starting with an effective pluricanonical divisor, we isolate a reduced
boundary component by the relative minimal model program.  Surface
adjunction supplies torsion on its reduction; local cohomology and
$p$-power gluing then give torsion on the full Cartier cycle.
A tame cover and primitive division produce $(V,D)$, and numerical
proportionality transfers the maximal-nef-dimension condition to $D$.
The intervening local surface arguments establish regularity in
codimension two and justify reduced-boundary adjunction.
The use of a connected primitive cycle follows the canonical-divisor
viewpoint of Miyaoka \cite[Definition~2.4 and Lemma~2.5]{miyaoka1988};
the torsion and covering arguments here take place in characteristic $p$.

\subsection{Resolution and field extension}

We use projective resolution in dimension three
\cite{cossartpiltant2008,cp2009,cp2019}.
The projective form that is unchanged over the regular locus is stated
explicitly in \cite[p.~1, introductory theorem]{cossart2008notes}.
It applies to our quasiprojective threefolds over the perfect field $k$.
For divisor supports on a regular threefold we use embedded resolution
\cite[Corollary~0.4]{cjs2009}: a reduced closed subscheme of dimension at
most two in a regular excellent scheme admits a projective modification,
unchanged off that subscheme, whose reduced inverse image is snc.
The resolutions may be chosen to be isomorphisms over the specified regular
open set.  In particular, pairs admit the log resolutions used below.

We also use the discrepancy criterion for a simple normal crossing
subpair on a smooth variety, allowing negative coefficients: coefficients
at most one give log canonicity, and coefficients strictly less than one
give the klt property.  One can see the relevant inequalities directly.
At the generic point of a center, use logarithmic differentials for the
components through it and regular differentials for the other coordinates.
At the generic point of a divisor on a smooth model, with uniformizer \(t\),
each pulled-back logarithmic differential is a regular differential plus
a multiple of \(dt/t\).  Their exterior product has at most a simple pole:
two occurrences of \(dt/t\) have zero exterior product.  The relative
Jacobian formula gives the discrepancy inequality.  Lowering a coefficient
adds its positive difference times the order of that component.
For an exceptional divisor whose center misses the boundary, the relative
Jacobian has positive order.  Indeed a birational morphism between smooth
models cannot be generically unramified along an exceptional divisor,
since it would then be quasi-finite there.  These observations also explain
why it suffices to compute thresholds on a log resolution.

We may replace \(k\) by an uncountable algebraically closed extension.
Here are the descent and preservation points needed for this reduction.
A normal variety over an algebraically closed field is geometrically
normal and geometrically integral.  In the present setting, geometric
normality also follows from Serre's criterion: the field is perfect,
residue fields of finite type are separably generated, and the relevant
flat base-change maps have regular fibers.  A smooth projective resolution
remains smooth and integral after extension.  The canonical rational line
bundle base-changes, since its restriction to the smooth open set does,
and reflexive extension determines it elsewhere.

On a resolution \(g:T\to X\), terminality gives
\[
 K_T=g^*K_X+\sum_E a_EE,\qquad a_E>0.
\]
The base-changed resolution has the same positive exceptional
discrepancies.  A further exceptional divisor has a positive relative
discrepancy over the smooth model and nonnegative contributions from
the effective divisor \(\sum_Ea_EE\); hence terminality is preserved.
Nefness is preserved, for example by writing a nef class as a limit of
ample rational classes, and intersection numbers preserve the numerical
dimension.  Finally, for every Cartier multiple \(qK_X\),
\[
 H^0(X_{k'},\cO_{X_{k'}}(qK_{X_{k'}}))
   =H^0(X,\cO_X(qK_X))\otimes_k k'.
\]
Thus pluricanonical sections descend in the required sense of equality
of dimensions.  If necessary, take a further multiple divisible by the
original Cartier index; the powers of two linearly independent sections
remain linearly independent, since their nonconstant ratio cannot acquire
a constant positive power in a function field with algebraically closed
constant field.  We do not assume that \(\Q\)-factoriality is preserved by
this extension: the modifications below supply it when needed.

\subsection{The maximal nef-dimension branch}

Nonvanishing \cite[Theorem~1.4]{xu2024} supplies an effective Cartier divisor
\[
 0\neq N_0\sim mK_X,\qquad m>0.
\]
It is nonzero because \(K_X\cdot H_X^2>0\) for an ample divisor \(H_X\).
It is nef, and \(N_0^2\cdot H_X=0\).
Every prime component \(T\) of \(N_0\) satisfies
\[
 N_0\cdot H_X\cdot T=0,
\]
because these nonnegative numbers have a positive weighted sum equal to
zero.  In fact \(N_0|_T\equiv0\).  To verify this for an integral curve
\(C\subset T\), choose a sufficiently large multiple of \(H_X|_T\) with
an effective Cartier member containing \(C\).  Its intersection with the
nef restriction is zero; all components contribute nonnegatively, so the
contribution of \(C\) is zero.

The nef reduction exists over our uncountable algebraically closed field
\cite[Definition~2.1 and Theorem~2.2]{xu2026}.
If \(n(X,K_X)\leq2\), the low-nef-dimension abundance theorem
\cite[Theorem~1.8]{xu2024} makes \(K_X\) semiample.  Since its numerical
dimension is one, some Cartier multiple has at least two sections.
We may therefore assume \(n(X,K_X)=3\).
The defining curve property of the nef reduction says in this case that
there is no \(N_0\)-trivial curve through a very general point of \(X\).
Here \emph{very general} means outside a countable union of proper closed
subsets.

\subsection{Isolating a boundary component}

We use the relative threefold MMP in characteristic greater than three,
including preservation of \(\Q\)-factorial dlt pairs for the boundary
driving the MMP, crepant dlt modifications, and the partial MMP with
nef restored boundary
\cite[Theorem~2.20, Proposition~2.21, Definition~2.23,
Lemma~2.24, and Theorem~2.27]{xu2026}.
The characteristic-five case is included through \cite{hw2022};
the corresponding results in characteristic greater than five include
\cite{hnt2020}.  Termination for the pseudoeffective lc divisor we use
is \cite[Theorem~1.5]{xu2024}.
The comparisons of pullbacks use the negativity lemma
\cite[Lemma~3.39]{kollarMori1998}.

We recall why the dlt preservation used here concerns negative steps.
On a common resolution of a negative divisorial contraction or flip,
the old log pullback minus the new one is effective by negativity.
On curves over the contraction target its degree is nonpositive, and
is strictly negative if the curve has a nonconstant image on a modified
side.  A component of a projective fiber not contained in this effective
divisor cannot meet it: its effective restriction is anti-nef, hence zero.
Connectedness of fibers then makes any fiber meeting the support lie
entirely in it.  Fibers over the modification locus do meet the support,
by the strict degree inequality.  Thus discrepancies strictly increase
for valuations centered in that locus.  A zero log discrepancy on the
output of a dlt step must consequently come from the unchanged open set,
where the required simple normal crossing behavior persists.
We will apply this to the lowered boundary.  For a larger boundary whose
log divisor is trivial along the steps we retain log canonicity by
crepancy, and later take a new dlt modification.

Let \(\lambda>0\) be the log canonical threshold of \(N_0\) on \(X\).
It is finite and rational, and is attained: on a log resolution it is
the minimum of finitely many positive rational discrepancy-to-order
ratios.  Thus \((X,\lambda N_0)\) is lc and is not klt.
Take a crepant projective \(\Q\)-factorial dlt modification
\[
 \mu:(W,B)\longrightarrow(X,\lambda N_0).
\]
All \(\mu\)-exceptional divisors occur in \(B\) with coefficient one.
One construction starts on a log resolution with the strict boundary
and all exceptional divisors reduced.  Its log divisor is the desired
pullback plus an effective exceptional discrepancy divisor.  The relative
dlt MMP removes that difference by negativity, leaving precisely the
crepant divisors.  This is the modification theorem cited above.

Put \(N=\mu^*N_0\) and \(\tau=m^{-1}+\lambda>0\).  Then
\begin{equation}
\label{eq:initial-dlt-pair}
 K_W+B\sim_{\Q}\tau N,\qquad \Supp B=\Supp N.
\end{equation}
Indeed an extracted log-canonical place must have center in \(\Supp N_0\),
since \(X\) is klt away from that divisor.  Its order on \(N_0\) is positive.
The modification is an isomorphism over the smooth open set
\(X_{\mathrm{reg}}\setminus\Supp N_0\): there it has no exceptional divisors,
and a projective small birational morphism to a locally factorial variety
is an isomorphism.  For the latter assertion, push down a relatively ample
Cartier divisor.  A multiple of its pushdown is Cartier, and its pullback
agrees with that multiple upstairs because there are no exceptional
divisors.  It therefore has degree zero on any contracted curve,
contradicting relative ampleness unless there are no such curves.

Since \((W,B)\) is dlt and not klt, choose a coefficient-one component \(S\).
Choose a rational \(\epsilon>0\) so small that both
\[
 B-\epsilon S\geq0,\qquad \tau N-\epsilon S\geq0.
\]
Run a \((K_W+B-\epsilon S)\)-MMP which is
\((K_W+B)\)-trivial.  In the partial-MMP notation the boundary is
\(B-\epsilon S\), and the restoration divisor is \(\epsilon S\).
The lowered pair is dlt and its log divisor is pseudoeffective by the
displayed effective representative.  Termination therefore applies.
There is no Mori fiber output, since a pseudoeffective divisor cannot be
anti-ample on the covering family of curves in a Mori fiber space.

We spell out what happens to \(N\).  On a common resolution of each step,
its two pullbacks agree.  Their difference is exceptional over the new
model.  A curve contracted over that model maps on the old model to a
curve over the step's contraction target, and has \(N\)-degree zero,
because the extremal ray is \(N\)-trivial.  The difference is therefore
numerically trivial over the new model, and is zero by negativity applied
to it and to its negative.  This proves the pullback assertion for flips
as well as divisorial contractions.

It follows that the transform of \(N\) remains nef and numerically trivial
on every component of its support.  For the latter assertion, lift a
curve on such a component to the common resolution.  Its image on the
old side is a point or a curve in the old support, where \(N\) is
numerically trivial.  The same pullback equality shows that the full
pair in \eqref{eq:initial-dlt-pair} remains crepant and lc.
The component \(S\) is not contracted: every contracted ray \(R\) satisfies
\[
 S\cdot R>0,
\]
whereas an exceptional divisor of a divisorial contraction has negative
degree on some contracted curve, again by negativity.
Moreover a negative curve for the effective divisor
\(\tau N-\epsilon S\) lies in its support.  Thus all steps are
isomorphisms off the support under consideration.

Write \(Y,B_Y,S_Y,N_Y\) for the output and its transforms.
The partial-MMP conclusion gives
\begin{equation}
\label{eq:isolated-nef-pair}
 K_Y+B_Y\sim_{\Q}\tau N_Y,\qquad
 K_Y+B_Y-\delta S_Y\ \text{nef}
\end{equation}
for some rational \(\delta>0\).
The variety \(Y\) is \(\Q\)-factorial and \((Y,B_Y)\) is lc.
If \(T\neq S_Y\) is a component of \(N_Y\), then
\(-S_Y|_T\) is nef by \eqref{eq:isolated-nef-pair} and \(N_Y|_T\equiv0\).
But \(S_Y|_T\) is effective and \(\Q\)-Cartier.  If nonzero, it has
positive degree against an ample divisor on the integral projective
surface \(T\), contradicting anti-nefness.  Consequently
\[
 S_Y\cap T=\varnothing\qquad(T\neq S_Y).
\]
The pair \((Y,S_Y)\) is lc, and, with \(s=\operatorname{mult}_{S_Y}N_Y>0\),
\begin{equation}
\label{eq:isolated-boundary}
 S_Y|_{S_Y}\equiv0,\qquad
 K_Y+S_Y\sim_{\Q}\tau s S_Y
 \quad\text{on a neighborhood of }S_Y .
\end{equation}

Take a new crepant projective \(\Q\)-factorial dlt modification
\[
 f:(Y_d,\Sigma)\longrightarrow(Y,S_Y),
\]
where \(\Sigma\) is the strict transform of \(S_Y\) together with the
exceptional divisors, all reduced.  We claim
\begin{equation}
\label{eq:second-dlt-support}
 \Supp f^*S_Y=\Supp\Sigma .
\end{equation}
For an extracted divisor \(E\), let \(a_E\) denote log discrepancy.
Log canonicity gives
\[
 0\leq a_E(Y,B_Y)
 =a_E(Y,S_Y)-\operatorname{ord}_E(B_Y-S_Y)
 =-\operatorname{ord}_E(B_Y-S_Y).
\]
Both terms on the right must vanish.
Crepancy with \((X,\lambda N_0)\) then gives
\(a_E(X,\lambda N_0)=0\); since \(X\) is klt,
\(\operatorname{ord}_E N_0>0\).
The pullback equalities just proved give
\(\operatorname{ord}_E N_Y>0\).
The other components of \(N_Y\) have the same support as \(B_Y-S_Y\)
and order zero at \(E\), so \(\operatorname{ord}_E S_Y>0\).
This proves \eqref{eq:second-dlt-support}.
The morphism \(f\) is an isomorphism off \(S_Y\), by the small-morphism
argument and \(\Q\)-factoriality of \(Y\).
Near \(\Sigma\) we therefore have
\begin{equation}
\label{eq:dlt-adjunction-input}
 K_{Y_d}+\Sigma\sim_{\Q}\tau s f^*S_Y,\qquad
 f^*S_Y|_\Sigma\equiv0 .
\end{equation}

\subsection{Local surface intersections and terminal regularity}

The adjunction calculation for \(\Sigma\) uses an intersection inequality
on excellent local surfaces, whose residue fields need not be perfect.
We prove that inequality first and also derive regularity of \(X\) in
codimension two.  For a resolution unchanged over
\(X_{\mathrm{reg}}\setminus\Supp N_0\), this regularity ensures that
exceptional divisors with centers outside \(\Supp N_0\) are disjoint
from its total pullback.

\begin{lemma}[Exceptional surface intersections]
\label{lem:exceptional-intersections}
Let \(P\) be an excellent normal local surface with closed residue field
\(l\), and let \(g:T\to P\) be a nontrivial projective regular resolution
which is an isomorphism away from the closed point.  Write \(G_i\) for its
exceptional integral curves, and take intersection degrees over \(l\).
Then \(G_i^2<0\) for every \(i\).  More generally, if an exceptional real
divisor \(Z=\sum_i z_iG_i\) has a positive coefficient, there is an index
\(j\) with \(z_j>0\) and \(Z\cdot G_j<0\).
\end{lemma}

\begin{proof}
The exceptional fiber is connected, by proper birationality and normality
of \(P\).  Choose a nonzero nonunit \(t\in\cO_P\) and write
\[
 \operatorname{div}_T(t)=A+\sum_i c_iG_i,\qquad c_i>0,
\]
where \(A\) is the strict-transform part.  It meets the exceptional fiber:
the strict transform of a component of the divisor of \(t\) has a point
over the closed point by properness.  For \(C=\sum_i c_iG_i\) we have
\[
 C\cdot G_i=-A\cdot G_i\leq0.
\]
Distinct curves have nonnegative intersections.  Connectedness therefore
shows \(G_i^2<0\): every vertex has either a neighbor or, if it is the only
vertex, a positive intersection with \(A\).

For the second assertion put \(a=\max_i(z_i/c_i)>0\), and let \(I\)
be the nonempty set where this maximum is attained.  Choose \(j\in I\)
which either meets \(A\) or has a neighbor outside \(I\).
Such a vertex exists by connectedness and the fact that \(A\) meets the
fiber.  Then
\[
 Z\cdot G_j
 =aC\cdot G_j+\sum_i(z_i-ac_i)G_i\cdot G_j<0.
\]
The term with \(i=j\) is zero; all other terms are nonpositive, and the
chosen vertex makes at least one term strictly negative.
\end{proof}

\begin{lemma}[Terminal threefolds are regular in codimension two]
\label{lem:terminal-codimension-two}
A terminal threefold in the present setting has only isolated
singularities.
\end{lemma}

\begin{proof}
Localize at a codimension-two point.  This gives an excellent normal local
surface \(P\), with residue field \(l\), and a projective regular resolution
\(g:T\to P\) obtained by localizing a threefold resolution.
The residue field \(l\) need not be perfect.  We keep all intersection
degrees over \(l\).

For an integral exceptional curve \(G\), divisor adjunction and curve
duality give
\begin{equation}
\label{eq:local-surface-adjunction}
 (K_T+G)\cdot G
 =2\dim_lH^1(G,\cO_G)-2\dim_lH^0(G,\cO_G).
\end{equation}
To specify the duality in this formula, use the dualizing complex on \(P\)
localized from the threefold and its exceptional inverse image on \(T\).
Transitivity for \(G\to T\to P\) and
\(G\to\Spec l\to P\) identifies the adjunction sheaf on \(G\) with
its dualizing sheaf over \(l\), up to tensoring with a one-dimensional
\(l\)-vector space.  This does not change degree.
Proper duality and the formula
\(\chi(\cL)-\chi(\cO_G)=\deg_l\cL\) prove
\eqref{eq:local-surface-adjunction}
\cite{hartshorne1966,hartshorne1977}.

Suppose \(K_T\cdot G<0\), and set
\(l'=H^0(G,\cO_G)\) and \(d=[l':l]\).
Both \(G^2\) and \(K_T\cdot G\) are negative multiples of \(d\).
Equation \eqref{eq:local-surface-adjunction} and
Lemma~\ref{lem:exceptional-intersections} force
\[
 H^1(G,\cO_G)=0,\qquad G^2=K_T\cdot G=-d.
\]
The line bundle \(\cO_G(-G)\) consequently has degree one over \(l'\).
Riemann--Roch gives at least two independent sections over \(l'\).
Every nonzero section has a zero scheme of length one.  Two independent
sections have no common zero: otherwise their effective Cartier divisors
coincide, so their ratio is a global unit, an element of \(l'^\times\).
They define a nonconstant map \(G\to\PP^1_{l'}\) of degree one.
It is finite, and a finite birational morphism to a normal curve is an
isomorphism.  Hence
\[
 G\simeq\PP^1_{l'},\qquad \cO_G(G)\simeq\cO_{\PP^1_{l'}}(-1).
\]

For completeness this curve can be contracted while retaining a
projective regular resolution over \(P\).  Choose a relatively very ample
Cartier divisor \(A\), put \(a=\deg_{l'}(A|_G)>0\), and set \(Q=A+aG\).
The divisor \(Q-G\) has positive degree on every component of the closed
fiber, and is relatively ample by the fiber criterion for ampleness
\cite{egaIII}.  The restriction \(Q|_G\) is trivial.
For every \(n\geq1\), a trivializing section of \(\cO_G(nQ)\) extends
successively to \(nG\), because the obstruction groups are
\[
 H^1\bigl(G,\cO_G(nQ-jG)\bigr)
   =H^1(\PP^1_{l'},\cO(j))=0\qquad(j\geq1).
\]
For large \(n\), relative Serre vanishing applied to \(n(Q-G)\) lifts
these sections from \(nG\) to \(T\).  They generate along \(G\).
Away from \(G\), the sections of \(n(Q-G)\), multiplied by the canonical
section of \(nG\), generate and separate points and tangent directions.
Thus \(nQ\) defines a projective contraction whose only nontrivial fiber
is \(G\).  Its Stein factorization has normal target and is an isomorphism
off \(G\).

Let \(q\) be the image of \(G\).  Formal functions identifies the completed
local ring at \(q\) with
\[
 \varprojlim_n H^0(nG,\cO_{nG}).
\]
The \(G\)-adic and inverse-image \(q\)-adic neighborhoods are cofinal,
since they have the same support.  The successive kernels have graded
pieces
\[
 H^0(G,\cO_G(-nG))=H^0(\PP^1_{l'},\cO(n)).
\]
There is no obstruction in \(H^1\), so the associated graded algebra is
the polynomial ring \(l'[x,y]\).
More explicitly, let \(F_n\) be the kernel of the map from the completed
local ring to \(H^0(nG,\cO_{nG})\).  Then \(F_1\) is its maximal ideal,
and the \(F_n\)-filtration is cofinal with the maximal-ideal-adic filtration.
Choose lifts of the two degree-one generators.
Generation of the graded algebra gives
\(F_1\subset(x,y)+F_n\) for every \(n\), by successive approximation,
using arbitrary lifts of coefficients from \(l'\).
Finitely generated ideals in the complete Noetherian local ring are
closed, so \(F_1=(x,y)\).  No embedding of the residue field \(l'\) into
the completed local ring is required.
Its dimension is two; hence
it, and then the local ring at \(q\), is regular.  This proves the claimed
contraction using formal functions and Serre vanishing
\cite[III, Sections~5 and~11]{hartshorne1977}.

There are only finitely many exceptional curves.  Contracting a
\(K_T\)-negative one whenever possible therefore yields a regular
resolution with \(K_T\) relatively nef.  If any exceptional curves remain,
terminality gives
\[
 E=K_T-g^*K_P=\sum_i a_iG_i,\qquad a_i>0.
\]
The discrepancies of surviving curves have not changed, since they are
read at their generic points.  Lemma~\ref{lem:exceptional-intersections}
gives \(E\cdot G_j<0\) for some \(j\), contradicting relative nefness.
There are no exceptional curves left, and the proper birational map is
an isomorphism.  Thus the codimension-two point was regular.
Normality already gives regularity in codimension one, so the remaining
singular locus of the projective threefold is finite.
\end{proof}

\subsection{Adjunction on the reduced boundary}

We return to the dlt pair \((Y_d,\Sigma)\) satisfying
\eqref{eq:dlt-adjunction-input}.  To apply surface abundance, we must
identify its restricted log canonical class with that of a projective
slc surface.

Let \(\sigma:\Sigma'\to\Sigma\) be the finite \(S_2\)-ification.
It is an isomorphism away from finitely many closed points
\cite[Proposition~2.14]{xu2026}.
Reduced-boundary adjunction gives
\begin{equation}
\label{eq:slc-adjunction}
 (K_{Y_d}+\Sigma)|_{\Sigma'}=K_{\Sigma'}+\Delta',
\end{equation}
where \((\Sigma',\Delta')\) is an slc surface pair
\cite[Lemma~2.19]{xu2026}.
Restrictions here and below mean rational line bundles, or sufficiently
divisible Cartier multiples.  We give the local verification needed
in this dlt situation.

Choose a log resolution and write
\begin{equation}
\label{eq:boundary-log-pullback}
 K_{\widetilde Y}+\widetilde\Sigma+\sum_i\beta_iE_i
   =h^*(K_{Y_d}+\Sigma),\qquad \beta_i\leq1.
\end{equation}
At a codimension-one point of \(\Sigma\), localize the ambient threefold
to a normal local surface \(P\), and write \(h:T\to P\) for the
localized regular resolution.  If the pair is not snc at this point,
dlt gives \(\beta_i<1\) on the exceptional curves over it.
Set
\[
 J=-\sum_i\lfloor\beta_i\rfloor E_i,\qquad
 M=J-\widetilde\Sigma,\qquad
 A_{\mathrm{fr}}=\sum_i\{\beta_i\}E_i.
\]
Here \(J\) is effective, and
\begin{equation}
\label{eq:surface-relative-vanishing-input}
 M-K_T-A_{\mathrm{fr}}\equiv_P0 .
\end{equation}

We claim \(R^1h_*\cO_T(M)=0\).
For every nonzero effective integral exceptional cycle \(Z\), the divisor
\(Z-A_{\mathrm{fr}}\) has a positive coefficient.
Lemma~\ref{lem:exceptional-intersections} gives a component \(E_j\) with
positive coefficient in \(Z\) such that
\((Z-A_{\mathrm{fr}})\cdot E_j<0\).
Using \eqref{eq:surface-relative-vanishing-input},
\[
 (M-Z+E_j)\cdot E_j>(K_T+E_j)\cdot E_j
   =\deg_l\omega_{E_j/l}.
\]
Curve duality gives
\(H^1(E_j,\cO_{E_j}(M-Z+E_j))=0\):
the dual space consists of sections of a line bundle of negative degree.
The exact sequence
\[
 0\longrightarrow\cO_{E_j}(M-Z+E_j)
 \longrightarrow\cO_Z(M)
 \longrightarrow\cO_{Z-E_j}(M)\longrightarrow0
\]
then proves \(H^1(Z,\cO_Z(M))=0\) by induction on the sum of
the coefficients of \(Z\).
Exceptional cycles give a cofinal system of neighborhoods of the
exceptional fiber.  Formal functions, followed by faithful completion
for the coherent higher direct image, proves the claim.

Normality gives \(h_*\cO_T(J)=\cO_P\): a rational function with possible
poles only along exceptional divisors is regular at every codimension-one
point of \(P\), hence regular on \(P\).
The preceding vanishing and
\[
 0\longrightarrow\cO_T(M)\longrightarrow\cO_T(J)
 \longrightarrow\cO_{\widetilde\Sigma}(J)\longrightarrow0
\]
give a surjection
\(\cO_P\to h_*\cO_{\widetilde\Sigma}(J)\).
It factors through \(h_*\cO_{\widetilde\Sigma}\), which is a subsheaf of
the target.  The strict boundary has no exceptional component, so its
proper map to the local boundary is quasi-finite and hence finite.
The kernel of
\(\cO_P\to h_*\cO_{\widetilde\Sigma}\) is precisely the ideal of the
reduced boundary.  The surjection thus identifies
\(\widetilde\Sigma\) with that boundary and also gives
\(\cO_{\widetilde\Sigma}=\cO_{\widetilde\Sigma}(J)\)
under their natural inclusion.
In particular \(J\) misses the strict boundary.

Two strict branches cannot cross above this non-snc point.
In a nontrivial regular surface resolution any such intersection over the
closed point lies on the exceptional fiber; the total divisor would then
have three components through a point on a surface, contrary to snc.
The trivial-resolution case would already be snc downstairs.
Thus the boundary is regular at a non-snc point of codimension one,
and every exceptional coefficient meeting its strict transform is
nonnegative.  At the snc points, the reduced boundary is nodal.
This proves that \(\Sigma'\) is nodal in codimension one.

Here is the log-line identification in \eqref{eq:slc-adjunction}.
Use rational residue forms to identify the plurilog-canonical line with
pluricanonical forms on each boundary component.
At snc points, Cartier divisor adjunction gives this identification,
including along the nodes.
At the other codimension-one points, restrict
\eqref{eq:boundary-log-pullback} to the smooth strict component.
The preceding local calculation identifies it with the boundary there
and shows that the resulting different is effective, since the exceptional
coefficients meeting it are nonnegative.
It has coefficient zero along the nodes.
The identification extends by the \(S_2\) property, using reflexive powers
with the indicated divisorial twist.  Its right side is therefore
\(\Q\)-Cartier with exactly the rational line bundle on the left side
of \eqref{eq:slc-adjunction}.

On the normalization of \(\Sigma'\), the boundary is the transform of
\(\Delta'\) together with the conductor, the latter with coefficient one.
This is the residue formula at a node, extended from codimension one.
Its log pullback to the smooth strict component
\(\widetilde\Sigma_i\) is
\[
 K_{\widetilde\Sigma_i}
 +\left((\widetilde\Sigma-\widetilde\Sigma_i)
                   +\sum_j\beta_jE_j\right)|_{\widetilde\Sigma_i}.
\]
This follows from smooth adjunction applied to
\eqref{eq:boundary-log-pullback}, with compatible rational residue forms,
so it is an equality of log pullbacks, not just numerical classes.
The displayed subpair is snc with every coefficient at most one.
The discrepancy criterion at the start of the section makes it lc.
Thus \(\Sigma'\) is \(S_2\), nodal in codimension one, and lc after
normalization with conductor; this proves the asserted slc property.

By \eqref{eq:dlt-adjunction-input}, the log canonical class in
\eqref{eq:slc-adjunction} is numerically trivial.
The surface \(\Sigma'\) is projective, since it is finite over \(\Sigma\).
Abundance for projective slc surfaces in positive characteristic
\cite[Theorem~3.1]{posva2024} makes this class semiample, hence
\(\Q\)-linearly trivial.  Indeed a globally generated numerically trivial
line bundle gives a constant morphism on each connected component and
is trivial there.
It follows from \eqref{eq:dlt-adjunction-input}, whose coefficient
\(\tau s\) is positive rational, that some positive Cartier multiple
\(L\) of \(f^*S_Y\) satisfies
\begin{equation}
\label{eq:reduced-boundary-triviality}
 \cO_{Y_d}(L)|_{\Sigma\setminus Z}\simeq\cO_{\Sigma\setminus Z}
\end{equation}
for a finite set \(Z\) of closed points.
Its support is \(\Sigma\), and \(L|_{\Sigma_i}\equiv0\) for every component.

\subsection{Torsion on the full Cartier cycle}

We now pass from \eqref{eq:reduced-boundary-triviality} to the actual
nonreduced divisor \(L\).

\begin{lemma}[Extending torsion across finitely many points]
\label{lem:full-cycle-torsion}
Let \(Y\) be a normal projective threefold over an algebraically closed
field of characteristic \(p>0\), and let \(L\) be an effective Cartier
divisor.  Suppose that \(L\) is numerically trivial on every integral
component of its support and that \(\cO_Y(L)|_{L_{\mathrm{red}}\setminus Z}\)
is trivial for a finite set of closed points \(Z\).
Then \(\cO_L(L)\) is torsion.
\end{lemma}

\begin{proof}
Replacing \(Z\) by \(Z\cap|L|\), we may assume \(Z\subset|L|\).
For \(z\in Z\), put \(R=\cO_{Y,z}\) and let \(\mathfrak m\) be its maximal
ideal.  Normality gives depth at least two, so
\[
 H^1_{\mathfrak m}(R)=0.
\]
Every nonmaximal localization of \(R\) is a normal local ring of dimension
at most two, and is consequently Cohen--Macaulay.
This implies that \(H^2_{\mathfrak m}(R)\) has finite length.
Explicitly, present \(R=A/I\), where \(A\) is a regular local ring of
dimension \(n\), essentially of finite type over the field.
The height of \(I\) is \(n-3\).
At a nonmaximal prime in \(\Spec R\), catenarity and
Auslander--Buchsbaum give projective dimension \(n-3\) over the
corresponding localization of \(A\).  Hence
\(\Ext_A^{n-2}(R,A)\) is supported only at the closed point.
It is a finite module, thus has finite length.
Local duality identifies \(H^2_{\mathfrak m}(R)\) with its Matlis dual,
up to completion, and proves the assertion
\cite{hartshorne1967}.

Let \(t_z\) be a local equation of \(L\).  Choose \(j\) so large that
\(t_z^{j-1}\) kills \(H^2_{\mathfrak m}(R)\) for every \(z\in Z\).
There are only finitely many points, so the same \(j\) works at all of
them.  The restriction sequences fit into a commutative diagram
\[
\begin{tikzcd}[column sep=small]
0 \arrow[r] &
\cO_Y((u-j)L) \arrow[r] \arrow[d] &
\cO_Y(uL) \arrow[r] \arrow[d,equal] &
\cO_{jL}(uL) \arrow[r] \arrow[d] & 0\\
0 \arrow[r] &
\cO_Y((u-1)L) \arrow[r] &
\cO_Y(uL) \arrow[r] &
\cO_L(uL) \arrow[r] & 0 .
\end{tikzcd}
\]
In local trivializations the left vertical map is multiplication by
\(t_z^{j-1}\).
The vanishing of the ambient \(H^1_Z\) makes the connecting maps from
the \(H^1_Z\) of the rightmost terms to the ambient \(H^2_Z\) injective.
The choice of \(j\) therefore gives
\begin{equation}
\label{eq:torsion-zero-comparison}
 H^1_Z(Y,\cO_{jL}(uL))
       \longrightarrow H^1_Z(Y,\cO_L(uL))
 \quad\text{equal to zero for every integer }u .
\end{equation}
The assertion is uniform in \(u\): twisting by \(uL\) only changes local
trivializations of these finite-length modules.

Fix this \(j\).
The kernel of
\(\cO_{jL\setminus Z}\to\cO_{L_{\mathrm{red}}\setminus Z}\)
is nilpotent.
Lift a chosen generator of the trivial line
\(\cO_Y(L)|_{L_{\mathrm{red}}\setminus Z}\) to local generating sections
on a finite open cover of \(jL\setminus Z\).
On overlaps their differences are in the nilpotent ideal.
For \(u=p^e\) sufficiently large, their \(u\)-th powers agree, since
the \(p^e\)-power map is additive in characteristic \(p\) and kills that
ideal.  The powers therefore give a trivializing section
\[
 s_j\in H^0(jL\setminus Z,\cO_{jL}(uL)).
\]
Project it to a section \(s\) on \(L\setminus Z\).
In the exact sequence of sections with support in \(Z\), its obstruction
to extending over \(L\) is the image of the obstruction for \(s_j\).
Equation \eqref{eq:torsion-zero-comparison} kills that image.
Thus \(s\) extends to \(H^0(L,\cO_L(uL))\).

This extension is a generator everywhere.
On each integral projective surface component \(T\) of \(L_{\mathrm{red}}\)
it is nonzero, since it is a generator off \(Z\).
If it vanished somewhere on \(T\), its zero scheme would be a nonempty
effective Cartier divisor: a nonzero section of an invertible sheaf
on an integral scheme is a non-zero-divisor in local trivializations.
Such a divisor has positive degree against an ample class, whereas
\(\cO_T(uL)\) is numerically trivial.  This is impossible.
No normality of \(T\) is needed for this intersection argument.
Nonvanishing after reduction implies invertibility on \(L\), by the
local criterion for a unit modulo a nilpotent ideal.
We conclude \(\cO_L(uL)\simeq\cO_L\), as required.
\end{proof}

Applying Lemma~\ref{lem:full-cycle-torsion} to \(Y_d\) and \(L\) proves
that \(\cO_L(L)\) is torsion.

\subsection{A tame cover and the primitive cycle}

Choose a smooth projective common resolution \(V_0\) of \(Y_d\) and \(X\).
It may be chosen to be an isomorphism over the common open set
\[
 X_{\mathrm{reg}}\setminus\Supp N_0
\]
and to resolve the relevant divisor supports.
Let \(\rho_0:V_0\to X\) be its map to \(X\).
The pullbacks of \(N_Y\) and \(N_0\) agree on \(V_0\), by the pullback
equalities through the MMP.

The section cutting out \(N_0\sim mK_X\) gives a rational section
\(\theta\) of \(\omega_{V_0}^{\otimes m}\) with
\begin{equation}
\label{eq:pluricanonical-section-divisor}
 \operatorname{div}(\theta)
     =\rho_0^*N_0+m(K_{V_0}-\rho_0^*K_X).
\end{equation}
Every exceptional center outside \(\Supp N_0\) is an isolated singular
point of \(X\), by Lemma~\ref{lem:terminal-codimension-two} and the choice
of the resolution.  The divisors over those centers are disjoint from
\(\rho_0^*N_0\).
Thus the support in \eqref{eq:pluricanonical-section-divisor} consists
of the total pullback support and possibly divisors disjoint from it.

Write
\[
 m=p^bm',\qquad b\geq0,\qquad (m',p)=1.
\]
Choose a nonzero rational generator \(\eta\) of
\(\omega_{V_0}^{\otimes p^b}\), so that
\(\theta=a\eta^{m'}\) for \(a\in k(V_0)^\times\).
Adjoin one root \(v\) of \(v^{m'}=a\), and normalize \(V_0\) in the
resulting field.  The extension is finite separable, possibly of degree
smaller than \(m'\).  Its normalization is projective.
Away from \(\Supp\operatorname{div}(\theta)\), this construction is an
integral component of the torsor of \(m'\)-th roots of a nowhere-vanishing
section of a line bundle.  It is finite étale there, since \(m'\) is
invertible in \(k\).

Resolve this normalization, and then the total divisor support, keeping
that regular étale open set unchanged.  Denote the resulting smooth
projective variety by \(V\), and its generically finite separable map to
\(V_0\) by \(\pi\).
The root \(v\eta\) is a rational section of
\(\pi^*\omega_{V_0}^{\otimes p^b}\).
The determinant of the differential of \(\pi\) is a nonzero rational
section of \(\omega_V\otimes\pi^*\omega_{V_0}^{-1}\);
it is nonzero because \(\pi\) is separable.
Multiplying \(v\eta\) by the \(p^b\)-th power of this Jacobian section produces a rational section
of \(\omega_V^{\otimes p^b}\).
Its divisor is supported over
\(\Supp\operatorname{div}(\theta)\), since \(\pi\) is étale off that
support.  In particular it is supported on the total pullback of \(N_0\),
together with divisors disjoint from it.

Let \(L_V\) be the pullback of \(L\) from \(Y_d\), and choose one connected
component \(C\) of \(\Supp L_V\).
The isolation of \(S_Y\) shows that \(C\) is disjoint from the remainder
of the total pullback support of \(N_0\).
Write the part of \(L_V\) supported on \(C\) as
\[
 aD,\qquad
 D=\sum_\alpha m_\alpha D_\alpha,\qquad
 a=\gcd_\alpha\bigl(\operatorname{mult}_{D_\alpha}L_V\bigr).
\]
Then \(D\) is a nonzero primitive effective Cartier divisor, since \(V\)
is smooth.  Its support is connected and snc by the final resolution.

The pullback \(L_V\) is numerically trivial on its support:
any curve there maps to a point or to a curve in \(\Sigma\), where
\(L\) is numerically trivial.  On a neighborhood of \(C\), \(L_V=aD\).
Consequently \(D|_{D_\alpha}\equiv0\).
For a curve not contained in \(\Supp D\), effectivity gives nonnegative
intersection with \(D\); for one contained in a component the degree is
zero.  Hence \(D\) is nef.  We have
\[
 D^2\cdot H=0,\qquad D\cdot H^2>0
\]
for every very ample divisor \(H\).

If \(\cO_L(tL)\simeq\cO_L\), its pullback is trivial on the Cartier
scheme \(L_V\).  Near \(C\), the inclusion \(D\subset aD=L_V\) is an
inclusion of closed subschemes.  Restricting the pulled-back
trivialization gives
\[
 \cO_D(taD)\simeq\cO_D.
\]
Thus \(\cO_D(D)\) is torsion.  This argument does not claim that its
order is unchanged under the cover or primitive division.

The rational canonical section constructed above has, on a neighborhood
of \(C\), divisor supported on \(D\); all other components of its divisor
are disjoint from \(C\).  It therefore gives
\eqref{eq:prepared-canonical}.  For any Cartier divisor \(A\) on \(V\),
\[
 p^bK_V\cdot D\cdot A
   =B_*\cdot D\cdot A
   =\sum_\alpha b_\alpha\,
       c_1(\cO_{D_\alpha}(D))\cdot A|_{D_\alpha}=0.
\]
This includes \(K_V^2\cdot D=0\).

\subsection{Numerical proportionality and the final interface}

Put \(\rho=\rho_0\circ\pi\) and \(N_{\mathrm{tot}}=\rho^*N_0\).
Both \(D\) and \(N_{\mathrm{tot}}\) are nef and nonzero.
The latter is numerically trivial on every component of its support:
every curve there maps to a point or to a curve in \(\Supp N_0\).
We have
\[
 D^2\cdot H=N_{\mathrm{tot}}^2\cdot H
           =D\cdot N_{\mathrm{tot}}\cdot H=0,
 \qquad D\cdot H^2>0,\quad N_{\mathrm{tot}}\cdot H^2>0.
\]
The surface Hodge index theorem implies that
\begin{equation}
\label{eq:prepared-proportionality}
 D\equiv cN_{\mathrm{tot}},\qquad
 c=\frac{D\cdot H^2}{N_{\mathrm{tot}}\cdot H^2}>0.
\end{equation}
Here is a way to check the numerical equivalence on every curve rather
than just on one general surface.
Given an integral curve \(C_0\subset V\), take \(a\gg0\) and an integral
surface \(T\in|aH|\) containing \(C_0\).
Such a surface exists by Bertini irreducibility with this prescribed
base locus \cite{jouanolou1983}: for sufficiently large \(a\), generators
of the twisted ideal of \(C_0\), multiplied by very ample sections,
give a linear system with no fixed divisor and with a birational map
to its image off \(C_0\).
The image has dimension three, so a general member is integral.

On a smooth projective resolution \(\widetilde T\) of \(T\), the pullbacks of \(D\) and
\(N_{\mathrm{tot}}\) are nef, isotropic, mutually orthogonal classes.
The pullback of \(H|_T\) has positive square.
The Hodge index form on \(N^1(\widetilde T)_{\R}\) has signature
\((1,\dim_{\R}N^1(\widetilde T)_{\R}-1)\)
\cite[V, Section~1]{hartshorne1977}; its restriction to the orthogonal
complement of any positive-square class is negative definite.
Two nonzero orthogonal isotropic classes with positive intersection
against that class are therefore positively proportional.
Their ratio is the number \(c\) in \eqref{eq:prepared-proportionality},
since intersection with the pullback of \(H|_T\) gives
\(aD\cdot H^2\) and \(aN_{\mathrm{tot}}\cdot H^2\).
Choose a curve on the resolution dominating \(C_0\); projection of
degrees proves \((D-cN_{\mathrm{tot}})\cdot C_0=0\).
This proves \eqref{eq:prepared-proportionality}.

Finally a very general point of \(V\) lies in the generically finite
locus of \(\rho\) and maps outside the countable union excluded by the
nef reduction on \(X\).
A curve through such a point has a curve as its image.
If it were \(D\)-trivial, \eqref{eq:prepared-proportionality} and the
projection formula would make its image \(N_0\)-trivial, a contradiction.
This establishes the remaining assertion of
Proposition~\ref{prop:prepared-cycle}.
The proof of that proposition is complete.
In the subsequent sections only $V,D,H$, the torsion order $r$, and the
local canonical relation \eqref{eq:prepared-canonical} are retained;
the auxiliary birational models are no longer needed.
